% Generated by ChimCanhCut Research Platform
% Template: IEEEtran — IEEE Transactions on Neural Networks and Learning Systems / TPAMI / TKDE
% Topic: Federated Learning under Concept Drift: Taxonomy, Challenges, and Future Directions

\documentclass[journal]{IEEEtran}

\usepackage[utf8]{inputenc}
\usepackage{cite}
\usepackage{amsmath,amssymb}
\usepackage{graphicx}
\usepackage{booktabs}
\usepackage{hyperref}
\usepackage{url}
\usepackage{multirow}
\usepackage{array}

\title{Federated Learning under Concept Drift: Taxonomy, Challenges, and Future Directions}
\author{[Author Names]}

\begin{document}
\maketitle

\begin{abstract}
% TODO: Enter your abstract via the Paper Metadata panel in the UI.
% Required: 150-250 words covering motivation, methodology, key findings, contributions.
\end{abstract}

\begin{IEEEkeywords}Federated Learning, Concept Drift, Non-IID Data, Distribution Shift\end{IEEEkeywords}

\section{Introduction}
\label{sec:introduction}

Federated Learning (FL) has emerged as a prominent distributed machine learning paradigm, enabling multiple clients to collaboratively train a global model while keeping raw data localized to preserve privacy. As FL deployment expands across large-scale applications, various structural paradigms have been developed to manage communication and aggregation. These architectural designs fundamentally dictate how information is routed and how models are updated across the network. Current FL architectures are categorized into four primary types:

*   **Centralized FL:** The standard topology where multiple clients update a single global model.
*   **Decentralized FL:** A peer-to-peer approach designed to enhance fault tolerance and reduce central bottlenecks.
*   **Hierarchical FL:** An architecture that organizes nodes into tiers, effectively reducing communication overhead and latency.
*   **Clustered FL:** A topology that groups clients based on data similarity to optimize model performance.

While these architectures address privacy and communication efficiency, traditional FL frameworks often assume that underlying data distributions remain stationary. In dynamic, real-world environments, this assumption rarely holds. Data streams generated by distributed edge devices are frequently subject to concept drift—a phenomenon where statistical properties change over time. In the context of FL, these distributional shifts can severely degrade the reliability of a global model. Concept drift manifests in four primary categories:

*   **Label drift:** Occurs when the prior probability of labels changes while the conditional distribution of features given the label remains constant.
*   **Feature drift:** Characterized by changes in the marginal distribution of input features.
*   **Real concept drift:** Involves a fundamental shift in the conditional distribution of the target given the input, altering the decision boundary mapping features to labels.
*   **Hybrid drift:** A complex, real-world scenario where multiple drift types occur simultaneously or sequentially, rendering previously learned mapping functions obsolete.

Mitigating these effects requires robust mechanisms for identifying shifts and updating models. However, the privacy-preserving nature of FL precludes direct inspection of client data, rendering traditional centralized drift-handling techniques largely inapplicable. Consequently, specialized strategies must be employed at the client or aggregation level. Detection strategies in FL are primarily categorized into statistical testing, performance-based monitoring, distance metrics, and implicit adaptation techniques. Once drift is identified—or as part of a proactive learning process—the system must employ an adaptation mechanism to maintain accuracy. The spectrum of adaptation mechanisms utilized in modern FL includes model retraining, ensemble methods, continual learning, knowledge distillation, and dynamic aggregation. Each approach offers distinct trade-offs regarding computational overhead, communication costs, and adaptation speed.

Despite the critical importance of addressing non-stationary data, a comprehensive systematization of the intersection between FL and concept drift is necessary to guide future research. To bridge this gap, this survey presents a taxonomy and critical review of FL under concept drift. By categorizing the problem space across four dimensions—drift type, detection strategy, adaptation mechanism, and FL architecture—this paper provides a structured framework for understanding current methodologies. Furthermore, this survey synthesizes prevailing challenges in adapting distributed models to shifting environments and outlines promising directions for developing robust, adaptive federated systems.

\section{Background and Related Work}
\label{sec:background}

**2. Background**

To systematically analyze the intersection of Federated Learning (FL) and concept drift, it is necessary to establish the foundational concepts governing distributed model training and dynamic data environments. This section delineates the core components of FL architectures, the theoretical categorization of concept drift, and the primary strategies utilized for drift detection and adaptation.

**2.1. Federated Learning Architectures**
Federated Learning is a distributed machine learning paradigm enabling multiple clients to collaboratively train a model without sharing raw, local data, thereby preserving privacy. The structural organization of participating nodes—the FL architecture—fundamentally impacts how data distributions and potential drifts are managed across the network. The primary architectures include:

*   **Centralized FL:** The standard topology in which a central server coordinates the learning process, aggregating model updates (e.g., gradients or weights) from all participating clients to refine a single global model.
*   **Decentralized FL:** A peer-to-peer framework that eliminates the central server. Clients communicate and aggregate updates directly with neighboring nodes, enhancing fault tolerance and mitigating central bottlenecks.
*   **Hierarchical FL:** An architecture introducing intermediate edge servers between clients and the central cloud. These edge servers perform partial, local aggregations before transmitting updates to the global server, effectively reducing communication overhead and latency.
*   **Clustered FL:** A topology that groups clients into distinct clusters based on similarities in local data distributions or learning objectives. Each cluster trains a specialized model, which is particularly beneficial in highly heterogeneous environments.

**2.2. Typology of Concept Drift**
In dynamic real-world environments, the underlying data distributions of FL clients are rarely static. Concept drift refers to the phenomenon where the statistical properties of the target variable or input features evolve over time. Understanding the specific drift type is essential for designing robust FL systems. Drift is generally categorized into four primary types:

*   **Label drift:** Occurs when the prior probability of class labels changes over time, while the conditional distribution of features given the label remains constant.
*   **Feature drift:** Characterized by changes in the distribution of input features, without altering the underlying decision boundary mapping features to labels.
*   **Real concept drift:** Involves a fundamental shift in the conditional distribution of the target given the features. This shift alters the actual decision boundary, rendering previously learned mapping functions obsolete.
*   **Hybrid drift:** A complex, real-world scenario where multiple types of drift (e.g., simultaneous feature and real concept drift) occur concurrently across the federated network.

**2.3. Detection Strategies**
Addressing concept drift in distributed environments first requires effective identification. The detection strategy employed within an FL framework dictates how shifts in data distribution are recognized across isolated clients. These strategies broadly encompass:

*   **Statistical testing:** Applies rigorous statistical measures to compare historical and current data distributions or analyzes shifts in model gradients to flag significant deviations.
*   **Performance-based:** Monitors the continuous degradation of the model's predictive accuracy, error rates, or loss metrics over time to infer the presence of drift.
*   **Distance metrics:** Calculates the mathematical distance between model weights, feature representations, or data summaries across different time steps to quantify distributional shifts.
*   **Implicit adaptation:** A strategy where the system does not explicitly flag drift but instead continuously updates the model to remain robust against evolving data.

\section{Problem Formulation}
\label{sec:problem_formulation}

To establish a rigorous foundation for analyzing federated learning (FL) in non-stationary environments, we must formally define the standard FL objective, incorporate the temporal dynamics of streaming data, and mathematically characterize the nature of concept drift.

In a standard stationary setting, an FL system comprises $N$ participating clients, where each client $i \in \{1, 2, \dots, N\}$ holds a local dataset $\mathcal{D}_i$ drawn from a local distribution $P_i(X, Y)$. Here, $X \in \mathcal{X}$ denotes the input feature space and $Y \in \mathcal{Y}$ the label space. The primary objective is to learn a global model, parameterized by $w$, that minimizes the empirical risk across all clients without centralizing the data:
$$ \min_w \sum_{i=1}^N \frac{|\mathcal{D}_i|}{|\mathcal{D}|} \mathcal{L}_i(w) $$
where $\mathcal{L}_i(w) = \mathbb{E}_{(x,y) \sim P_i}[\ell(w; x, y)]$ is the expected loss for client $i$, and $|\mathcal{D}|$ is the total number of samples. This objective serves as the basis for various structural topologies, classified in our taxonomy as the *FL Architecture*:
*   **Centralized FL:** The standard topology where clients update a single global model.
*   **Decentralized FL:** A peer-to-peer architecture designed to improve fault tolerance and reduce central bottlenecks.
*   **Hierarchical FL:** An architecture that leverages multi-tier aggregation to mitigate communication overhead and latency.
*   **Clustered FL:** A topology that groups clients based on data similarity to improve convergence.

In real-world dynamic environments, the assumption that $P_i(X, Y)$ remains static is frequently violated. To formulate concept drift, we introduce a temporal dimension $t$. At any time $t$, client $i$ receives a stream of data drawn from a time-varying joint distribution $P_{i,t}(X, Y)$. Concept drift is formally defined as a change in this joint probability distribution, such that for a client $i$ at two distinct time points $t$ and $t+\Delta$:
$$ P_{i,t}(X, Y) \neq P_{i,t+\Delta}(X, Y) $$

Based on the decomposition of the joint distribution—either $P(X, Y) = P(Y|X)P(X)$ or $P(X, Y) = P(X|Y)P(Y)$—we categorize *Drift Type* into four primary classes:
*   **Label drift:** Occurs when the prior probability of the target labels changes, $P_{i,t}(Y) \neq P_{i,t+\Delta}(Y)$, while the class-conditional probability remains constant.
*   **Feature drift:** Characterized by changes in the marginal distribution of input features, $P_{i,t}(X) \neq P_{i,t+\Delta}(X)$, while the conditional distribution mapping features to labels remains constant.
*   **Real concept drift:** Involves a fundamental shift in the conditional distribution, $P_{i,t}(Y|X) \neq P_{i,t+\Delta}(Y|X)$, rendering previously learned mapping functions obsolete.
*   **Hybrid drift:** A complex, real-world scenario where multiple drift types occur concurrently within the federated network.

To maintain model utility, the FL system must continuously solve a dynamic optimization problem, necessitating two critical operational phases: detection and adaptation.

Let $\Phi$ represent a *Detection Strategy* function. The goal of $\Phi$ is to monitor local or global states and trigger a drift signal $\delta_{i,t} \in \{0, 1\}$. The formulation of $\Phi$ varies by strategy: it may evaluate distributional discrepancies via statistical testing, monitor performance degradation, quantify shifts using distance metrics, or bypass explicit signaling through implicit adaptation.

Upon detection ($\delta_{i,t} = 1$), the system invokes an *Adaptation Mechanism*, denoted as $\Psi$, to update the model parameters $w_t \rightarrow w_{t+1}$ such that the risk on the new distribution $P_{i, t+\Delta}$ is minimized. The formulation of $\Psi$ dictates how the network recovers from the shift, depending on the severity of the change.

\section{Methodology}
\label{sec:methodology}

\# Methodology

This systematic literature review (SLR) adheres to the Preferred Reporting Items for Systematic Reviews and Meta-Analyses (PRISMA) 2020 guidelines. Our methodology ensures transparency, reproducibility, and rigor throughout the selection and synthesis of literature concerning Federated Learning (FL) under concept drift.

\#\# 1. Search Strategy
To achieve comprehensive domain coverage, we conducted systematic searches across three academic repositories: Semantic Scholar, arXiv, and IEEE Xplore. The search window spanned January 2015 to January 2026, capturing the evolution of FL from its inception to contemporary state-of-the-art developments.

We utilized a combination of Boolean operators and keyword clusters to target the intersection of distributed machine learning and non-stationary data distributions. The primary query structure was:
*("Federated Learning" OR "Distributed Learning") AND ("Concept Drift" OR "Data Drift" OR "Non-stationary" OR "Distribution Shift" OR "Temporal Dynamics")*

We executed 30 distinct search queries, iterating through terminological variations to mitigate the risk of omitting relevant literature due to inconsistent field nomenclature.

\#\# 2. PRISMA Flow and Study Selection
The selection process followed the four-stage PRISMA flow diagram. Initially, 191 records were identified.

*   **Identification:** 191 records were retrieved; no duplicates were identified.
*   **Screening:** Titles and abstracts were screened to exclude non-English papers, those lacking accessible abstracts, or publications predating 2015.
*   **Eligibility:** We performed full-text assessments on the remaining records to ensure alignment with our research objectives.
*   **Included:** Studies were categorized into three tiers: *direct* (explicitly addressing concept drift in FL), *adjacent* (addressing related challenges like client heterogeneity or catastrophic forgetting), and *foundational* (providing theoretical bases for FL or drift detection).

\#\# 3. Inclusion and Exclusion Criteria
To maintain the survey's focus and quality, we applied the following criteria:

**Inclusion Criteria:**
*   **Peer-reviewed or arXiv preprints:** Balancing established research with rapid-release, cutting-edge findings.
*   **Relevance:** Studies must contribute to understanding FL under non-stationary conditions via algorithmic innovation, theoretical analysis, or empirical benchmarking.
*   **Temporal Scope:** Published between 2015 and 2026.
*   **Accessibility:** Full abstracts must be available for evaluation.

**Exclusion Criteria:**
*   **Duplicates:** To prevent synthesis bias.
*   **Non-English:** To ensure accurate technical interpretation.
*   **Pre-2015:** To maintain relevance to the current FL landscape.
*   **Inaccessible Abstracts:** Papers with unverifiable methodologies were excluded to preserve review integrity.

\#\# 4. Quality Assessment
We evaluated the quality of included studies using a multi-dimensional rubric based on:
1.  **Methodological Rigor:** The clarity of the proposed FL architecture and the robustness of drift detection or adaptation mechanisms.
2.  **Experimental Validation:** The utilization of standard datasets (e.g., FEMNIST, CIFAR-10/100, or synthetic drift benchmarks) and appropriate baselines.
3.  **Contribution:** The extent to which the research addresses specific concept drift challenges, such as communication overhead, privacy preservation, or model convergence stability.

\begin{figure}[!t]
\centering
\includegraphics[width=0.85\columnwidth]{figures/prisma_flow}
\caption{PRISMA 2020 flow diagram of the systematic literature search.}
\label{fig:prisma}
\end{figure}

\section{Taxonomy}
\label{sec:taxonomy}

To systematically analyze the intersection of federated learning (FL) and concept drift, we propose a comprehensive taxonomy that categorizes the problem space along four primary dimensions: Drift Type, Detection Strategy, Adaptation Mechanism, and FL Architecture. This multi-dimensional framework provides a structured lens through which to evaluate how non-stationary data distributions are managed across distributed, privacy-preserving networks.

**Drift Type**
The first dimension characterizes the fundamental statistical nature of the distributional shift occurring within or across FL clients. We classify these shifts into four distinct categories:
*   **Label drift:** Refers to variations in the prior probability of classes or target variables, often manifesting as dynamic class imbalance across different clients over time.
*   **Feature drift:** Occurs when the distribution of the input features changes, while the underlying mapping between features and labels remains constant.
*   **Real concept drift:** Represents a fundamental shift in the conditional distribution of the target variable given the input features, meaning the actual decision boundary has evolved and historical mappings are no longer valid.
*   **Hybrid drift:** Encompasses scenarios where multiple forms of drift, such as simultaneous feature and real concept drift, occur concurrently, presenting the most complex challenge for collaborative model training.

**Detection Strategy**
The second dimension focuses on the strategies employed to identify the onset, location, and severity of concept drift in a federated environment:
*   **Statistical testing:** Relies on mathematical tests comparing historical and current data distributions—often approximated via model gradients or local updates due to privacy constraints—to flag significant statistical deviations.
*   **Performance-based detection:** Monitors the degradation of local or global model efficacy metrics, such as accuracy or loss, triggering an alert when performance drops below a predefined threshold.
*   **Distance metrics:** Quantify the divergence between client updates, weight vectors, or data representations in latent spaces to identify drifting clients.
*   **Implicit adaptation:** Strategies that do not explicitly flag discrete drift events; instead, they utilize continuous optimization techniques that inherently absorb and mitigate the impact of shifting distributions without requiring a formal detection trigger.

**Adaptation Mechanism**
Once drift is detected or anticipated, the system must employ an adaptation mechanism, which forms the third dimension of our taxonomy:
*   **Model retraining:** A direct approach involving the partial or complete retraining of global or local models using newly acquired, representative data.
*   **Ensemble methods:** Maintain multiple models or sub-models across the network, dynamically adjusting their aggregation weights or voting power based on their relevance to the current data distribution.
*   **Continual learning:** Techniques integrated to allow models to learn new concepts sequentially, employing regularization or memory buffers to prevent the catastrophic forgetting of previously acquired knowledge.
*   **Knowledge distillation:** Facilitates adaptation by transferring knowledge from historically robust models to newer models adapting to current data, ensuring stability.
*   **Dynamic aggregation:** Modifies server-side federated aggregation protocols, adjusting client weighting mechanisms to prioritize updates from clients that reflect the most current and relevant conceptual state.

**FL Architecture**
The final dimension delineates the structural topology of the federated network, which fundamentally dictates how drift detection and adaptation are orchestrated:
*   **Centralized FL:** The standard topology where a central server coordinates updates, aggregates models, and manages global drift adaptation strategies.
*   **Decentralized FL:** A peer-to-peer architecture that removes the central server, allowing clients to update a single global model through direct communication, thereby enhancing fault tolerance and reducing central bottlenecks.
*   **Hierarchical FL:** An architecture that organizes clients into tiers to manage communication overhead and latency.
*   **Clustered FL:** A topology that groups clients with similar data distributions or drift patterns, allowing for specialized model training and more targeted adaptation.

\begin{table}[!t]
\renewcommand{\arraystretch}{1.3}
\caption{Proposed Multi-Dimensional Taxonomy}
\label{tab:taxonomy}
\centering
\begin{tabular}{p{3cm}p{9cm}}
\toprule
\textbf{Dimension} & \textbf{Categories} \\
\midrule
\textbf{Drift Type} & Label Drift, Feature Drift, Real Concept Drift, Hybrid Drift \\
\textbf{Detection Strategy} & Statistical Testing, Performance-based, Distance Metrics, Implicit Adaptation \\
\textbf{Adaptation Mechanism} & Model Retraining, Ensemble Methods, Continual Learning, Knowledge Distillation, Dynamic Aggregation \\
\textbf{FL Architecture} & Centralized FL, Decentralized FL, Hierarchical FL, Clustered FL \\
\bottomrule
\end{tabular}
\end{table}

\section{Literature Review}
\label{sec:literature_review}

The literature concerning federated learning (FL) under concept drift has expanded significantly as researchers seek to deploy distributed machine learning models in dynamic, non-stationary environments. To systematically evaluate the current state of research, existing studies can be analyzed through a multi-dimensional taxonomy comprising four core aspects: the type of drift, the detection strategy employed, the adaptation mechanism utilized, and the underlying FL architecture.

**Drift Type**
In non-stationary distributed environments, the literature identifies four primary variations of concept drift. First, label drift occurs when the marginal distribution of target variables changes across clients or over time, a common scenario in highly heterogeneous FL networks. Second, feature drift involves shifts in the input data distribution while the underlying relationship between features and labels remains constant. Third, real concept drift is characterized by a fundamental change in the conditional distribution of the target given the input features; this shift invalidates previously learned decision boundaries and necessitates substantial model updates. Finally, hybrid drift encompasses the simultaneous occurrence of the aforementioned phenomena, representing the most complex and realistic scenarios addressed in recent literature.

**Detection Strategy**
A critical component of managing concept drift in FL is the timely and accurate identification of distributional shifts across decentralized nodes. The literature reveals four primary detection strategies. Statistical testing methods are frequently employed to identify significant deviations in data distributions at the client level before global aggregation occurs. Alternatively, performance-based detection relies on monitoring the degradation of model metrics—such as accuracy or loss—on local validation sets, triggering adaptation protocols when predefined thresholds are breached. Distance metrics are also utilized to quantify the divergence between historical and current data representations or model weights. In contrast to explicit detection triggers, a subset of studies advocates for implicit adaptation, where the learning algorithm is inherently designed to adjust continuously to new data distributions without requiring a discrete drift detection mechanism.

**Adaptation Mechanism**
Once drift is detected or anticipated, FL systems must employ robust mechanisms to maintain model efficacy. The literature categorizes these into several distinct approaches. Model retraining is the most direct method, involving the partial or complete retraining of the global model using newly acquired client data. To mitigate catastrophic forgetting and handle recurring concepts, ensemble methods are widely discussed; these approaches maintain multiple models or sub-models and aggregate their predictions based on the current environmental context. Continual learning techniques are increasingly integrated into FL frameworks to allow models to sequentially learn new concepts while preserving previously acquired knowledge. Furthermore, knowledge distillation is leveraged to transfer information from historical, stable models to current ones, ensuring stability during the adaptation phase. Finally, dynamic aggregation strategies modify the server-side aggregation process; by adjusting client weights based on their relevance to the new concept, the central server can prioritize updates from clients experiencing the most representative data shifts.

**FL Architecture**
The structural design of an FL system significantly influences both its vulnerability to concept drift and the strategies available for mitigation. The primary architectures include:

*   **Centralized FL:** The standard topology where a central server coordinates the aggregation of updates from clients to update a single global model.
*   **Decentralized FL:** A peer-to-peer approach that eliminates the central server to improve fault tolerance and reduce central bottlenecks.
*   **Hierarchical FL:** An architecture that leverages multi-tier aggregation to reduce communication overhead and latency.
*   **Clustered FL:** A topology that groups clients with similar data distributions to improve model personalization and adaptation efficiency.

\section{Critical Analysis}
\label{sec:critical_analysis}

A critical examination of the intersection between federated learning (FL) and concept drift reveals significant structural trade-offs across the dimensions of the established taxonomy: Drift Type, Detection Strategy, Adaptation Mechanism, and FL Architecture. While existing frameworks attempt to address non-stationary environments, the isolated application of these taxonomic dimensions often leads to suboptimal system performance and exposes fundamental limitations in current methodologies.

Analyzing the relationship between Drift Types and Detection Strategies exposes a primary diagnostic bottleneck. Performance-based detection strategies are frequently utilized due to their operational simplicity; however, they inherently lack the granularity required to distinguish between Real Concept Drift and mere Label Drift. When a global model's performance degrades, performance-based metrics alone cannot identify whether the underlying conditional distribution has shifted or if local label proportions have merely skewed. Conversely, Statistical Testing and Distance Metrics offer a more rigorous mathematical foundation for identifying Feature Drift and complex Hybrid Drift. Yet, computing these metrics across distributed clients without violating the strict privacy constraints of FL remains a critical tension point. Implicit Adaptation bypasses explicit detection entirely, but this strategy often results in catastrophic forgetting or delayed convergence when the network is subjected to severe Real Concept Drift.

Furthermore, the taxonomy of Adaptation Mechanisms highlights a pervasive trade-off between computational overhead and adaptation agility. Model Retraining, while straightforward, is highly inefficient in dynamic environments characterized by frequent Hybrid Drift, as it discards previously acquired knowledge and consumes excessive communication bandwidth. Ensemble Methods provide robustness against recurring concepts, but they impose severe memory and computational burdens on resource-constrained edge devices. Continual Learning and Knowledge Distillation offer promising theoretical pathways to mitigate catastrophic forgetting during Real Concept Drift; however, their efficacy is heavily dependent on the availability of representative anchor data, which is fundamentally restricted in privacy-preserving FL settings. Dynamic Aggregation emerges as a highly synergistic mechanism, particularly when paired with Distance Metrics, as it allows the server to re-weight client updates based on their localized drift status. Nevertheless, Dynamic Aggregation struggles to maintain global convergence under conditions of rapid, uncoordinated Hybrid Drift across disparate client distributions.

The chosen FL Architecture fundamentally dictates the feasibility and efficiency of the aforementioned detection and adaptation strategies. Centralized FL remains the standard paradigm, yet it presents communication bottlenecks when coordinating global Statistical Testing for drift detection. Decentralized FL alleviates this bottleneck but severely complicates the consensus required for Ensemble Methods and Dynamic Aggregation, as the absence of a central orchestrator makes it difficult to align drifting concepts across peer-to-peer networks. 

Hierarchical FL and Clustered FL present the most structurally sound architectures for managing complex drift scenarios. Clustered FL, in particular, demonstrates a natural affinity for handling Hybrid Drift. By grouping clients into clusters based on Distance Metrics, the system can isolate drifting clients and apply localized Model Retraining or Knowledge Distillation without perturbing the global model or negatively impacting stationary clients. Similarly, Hierarchical FL localizes drift detection to intermediate edge servers, thereby reducing the latency of Adaptation Mechanisms and enabling more responsive Statistical Testing. 

Ultimately, this critical analysis indicates that current methodologies often treat concept drift in FL as a modular problem rather than a systemic one. The reliance on Implicit Adaptation within Centralized FL architectures fails to adequately address the nuances of Real Concept Drift and Feature Drift. The development of robust systems requires moving toward tightly coupled frameworks where the FL Architecture (such as Clustered FL) natively supports the Detection Strategy to trigger the most appropriate Adaptation Mechanism tailored to the specific Drift Type.

\section{Future Directions and Open Problems}
\label{sec:future_directions}

While the proposed taxonomy categorizes the current landscape of Federated Learning (FL) under concept drift into specific drift types, detection strategies, adaptation mechanisms, and architectural designs, several critical research gaps remain. The following future directions highlight key opportunities for advancing the robustness, efficiency, and scalability of FL systems operating in non-stationary environments.

**Tackling Hybrid and Real Concept Drift**
Current methodologies frequently isolate specific drift modalities, focusing predominantly on independent Label Drift or Feature Drift. However, real-world federated networks frequently exhibit Hybrid Drift, where input feature distributions and underlying decision boundaries shift simultaneously across different nodes. Future research must prioritize the development of unified frameworks capable of disentangling the individual components of Hybrid Drift. Furthermore, addressing Real Concept Drift—where the fundamental relationship between features and labels changes—remains a significant challenge in distributed settings. Future investigations should focus on distinguishing Real Concept Drift from transient data noise or benign statistical heterogeneity, ensuring that federated models do not erroneously adapt to localized anomalies while ignoring genuine systemic shifts.

**Advancing Proactive and Privacy-Preserving Detection Strategies**
The taxonomy of Detection Strategies highlights a current reliance on Performance-based metrics and Statistical Testing. A major future direction involves transitioning from reactive Performance-based detection, which inherently suffers from adaptation latency and performance degradation, toward proactive Implicit Adaptation. Implicit Adaptation mechanisms that seamlessly integrate drift detection directly into the model optimization process—without requiring explicit trigger thresholds—represent a highly promising frontier. Additionally, while Distance Metrics are valuable for quantifying distribution shifts between clients and the server, future work must ensure these metrics operate exclusively on model parameters or gradients to strictly preserve data privacy. Developing lightweight, privacy-preserving Distance Metrics that function efficiently without exposing raw client data is essential for the next generation of FL systems.

**Synergistic and Resource-Efficient Adaptation Mechanisms**
In terms of Adaptation Mechanisms, traditional Model Retraining imposes severe computational and communication overheads, making it impractical for edge devices. Future research should explore the intersection of Continual Learning and Knowledge Distillation to formulate more efficient adaptation strategies. By leveraging Knowledge Distillation, FL systems can transfer updated, drift-adapted representations to clients without the need for full-scale Model Retraining. Simultaneously, Continual Learning techniques must be integrated to mitigate catastrophic forgetting when the global model adapts to Real Concept Drift. Moreover, combining Ensemble Methods with Dynamic Aggregation offers a robust pathway for creating resilient systems. Future protocols could utilize Dynamic Aggregation to intelligently weight ensemble components based on real-time drift severity, allowing clients to maintain personalized, drift-resistant models.

**Architectural Innovations for Localized Drift Response**
The underlying FL Architecture plays a pivotal role in how concept drift is managed and mitigated. While Centralized FL provides a straightforward aggregation mechanism, it often creates bottlenecks in detecting and responding to heterogeneous drift across thousands of clients. Future directions must increasingly leverage Decentralized FL and Hierarchical FL architectures to enable rapid, localized responses to non-stationary data. Specifically, Clustered FL emerges as a critical area for future exploration. By dynamically grouping clients into clusters based on shared drift trajectories, Clustered FL can isolate the effects of localized concept drift. Future research should focus on algorithms that dynamically reassign clients across clusters in real-time as their local data distributions undergo Feature Drift or Label Drift, thereby maintaining high local performance without degrading the broader federated network.

\section{Conclusion}
\label{sec:conclusion}

In this survey, we have provided a comprehensive overview of the intersection between Federated Learning (FL) and concept drift—a critical challenge arising when the statistical properties of target variables or input data evolve over time. While FL has emerged as a leading paradigm for privacy-preserving, distributed machine learning, its foundational assumption of static data distributions is frequently violated in real-world applications. To systematically map how the research community addresses these dynamic environments, this paper introduced a multi-dimensional taxonomy encompassing drift types, detection strategies, adaptation mechanisms, and FL architectures.

Our proposed taxonomy first categorizes the fundamental nature of non-stationary environments in distributed networks. We identified four primary types of drift: label drift, feature drift, real concept drift, and hybrid drift. Discerning the specific nature of the drift is essential, as the shift in input distribution, target distribution, or the conditional relationship between them dictates the required algorithmic response. Furthermore, this survey highlights that the impact and management of these drifts are heavily influenced by the underlying FL architecture. The literature demonstrates that concept drift manifests and must be mitigated differently depending on whether the system employs:
*   **Centralized FL:** The standard topology where clients update a single global model.
*   **Decentralized FL:** A peer-to-peer approach designed to enhance fault tolerance and reduce central bottlenecks.
*   **Hierarchical FL:** An architecture that organizes nodes into layers to reduce communication overhead and latency.
*   **Clustered FL:** A topology that groups clients with similar data distributions, which can inherently isolate localized drifts, whereas centralized and decentralized architectures require distinct synchronization strategies to propagate drift awareness across the broader network.

A critical component of managing non-stationary distributed data is the timely and accurate identification of distribution shifts. The survey detailed four primary detection strategies currently utilized in the field. Statistical testing and distance metrics provide rigorous mathematical frameworks for comparing historical and current data distributions, though they often require careful design to comply with FL's strict privacy constraints. Alternatively, performance-based detection relies on monitoring the degradation of model metrics—such as accuracy or loss—as a proxy for underlying data shifts. Finally, implicit adaptation strategies bypass explicit detection, embedding drift resilience directly into the optimization process so that global or local models continuously adjust to new concepts without requiring a formal trigger.

Once drift is detected or anticipated, the FL system must employ robust adaptation mechanisms to maintain model efficacy. Our taxonomy categorizes these approaches into model retraining, ensemble methods, continual learning, knowledge distillation, and dynamic aggregation. Model retraining remains a foundational approach, though it can be communication-intensive across distributed nodes. Ensemble methods offer flexibility by maintaining multiple models to represent different concepts, allowing clients to weight outputs based on the currently active concept. Continual learning techniques are increasingly integrated into FL to learn new concepts sequentially while mitigating the catastrophic forgetting of previous knowledge. Knowledge distillation provides a communication-efficient means to transfer updated concept representations between clients and servers. Lastly, dynamic aggregation adjusts the server-side weighting of client updates, prioritizing contributions from clients that have successfully adapted to the new concept while marginalizing obsolete updates.

In conclusion, ensuring the robustness of Federated Learning under concept drift is paramount for its successful deployment in dynamic, real-world ecosystems. By structuring the existing literature through the lenses of drift type, detection, and adaptation, this survey provides a roadmap for future research in resilient distributed intelligence.

\bibliographystyle{IEEEtran}
\bibliography{references}

\end{document}